\documentclass[11pt,a4paper]{article}

% Idioma e fontes (compilar com XeLaTeX)
\usepackage{fontspec}
\usepackage[brazilian]{babel}

% Matemática, figuras e links
\usepackage{amsmath, amssymb}
\usepackage{graphicx}
\usepackage[hidelinks]{hyperref}

\title{Título do artigo}
\author{Nome Sobrenome}
\date{\today}

\begin{document}

\maketitle

\begin{abstract}
  Resumo breve do trabalho: o problema, a abordagem e os resultados principais.
\end{abstract}

\section{Introdução}
\label{sec:introducao}

Este é um artigo de exemplo. Os acentos, o ç e o til (ã, õ) funcionam
diretamente graças ao XeLaTeX. Para mais detalhes, veja a seção~\ref{sec:desenvolvimento} e a
referência~\cite{knuth1984}.

\section{Desenvolvimento}
\label{sec:desenvolvimento}

Uma equação numerada:
\begin{equation}
  \label{eq:euler}
  e^{i\pi} + 1 = 0
\end{equation}

e um sistema alinhado:
\begin{align}
  \nabla \cdot \vec{E} &= \frac{\rho}{\varepsilon_0}, \\
  \nabla \cdot \vec{B} &= 0.
\end{align}

A equação~\eqref{eq:euler} relaciona cinco constantes fundamentais.

\section{Conclusões}

Resuma aqui as conclusões principais.

\bibliographystyle{plain}
\bibliography{referencias}

\end{document}
